\section{Related work}
\label{sec:related-work}

Structural and longitudinal causal frameworks provide the counterfactual language for interventions and treatment regimes \citep{Pearl2009,Robins2000,vanderLaan2003,vanderLaan2011}. Within that language, theoretical work develops identification and estimation for stochastic and modified policies in continuous-treatment settings, including rules whose assigned dose depends on the observed natural value \citep{Diaz2012,Haneuse2013,Young2014,DiazWilliamsHoffmanSchenck2023}. Software and applied guidance make these policies operational and illustrate threshold interventions motivated by positivity and support constraints \citep{WilliamsDiaz2023lmtp,HoffmanSalazarBarretoWilliamsRudolphDiaz2024}. The present paper studies the \leanref{S-1}{observed-data experiment generated by a lower-threshold modified treatment policy} in which the deterministic rule clamps natural treatment values below the threshold to the threshold itself. That exact clamp creates a point mass at the boundary of the post-policy distribution, so the estimand combines an ordinary retained natural-course component with a boundary regression value multiplied by the mass collapsed by the policy.

A second line of work studies continuous-treatment dose-response functionals and the role of overlap in semiparametric and nonparametric inference \citep{Imbens2000,Hirano2005,Imai2004,KennedyMaMcHughSmall2017,BonviniKennedy2026}. Recent natural-value-policy work develops root-\(n\) targeted-learning inference under its regularity and nuisance-rate conditions \citep{Diaz2026}, while positivity-violation work separates an in-support point-identified component from a Lipschitz-bounded extrapolation component \citep{Koo2026}. Threshold-response methods target support-restricted stochastic functionals with efficient estimation and simultaneous bands \citep{vanderLaanZhangGilbert2023}, and stochastic-shift work targets effect modification with valid confidence intervals \citep{McCoyZhangHubbardvanderLaanSchuler2026}. The present estimand has a different post-policy law: its deterministic clamp maps the entire lower tail to one moving boundary point, producing a Dirac component whose weight and regression value must be learned together. Under a global two-sided polynomial envelope, the delivered results characterize minimax risk and uniformly honest expected length for this unsmoothed clamp functional.

The nonparametric component is closely related to regression with degenerate design. For smoothness \(s\) and a design density regularly varying with exponent \(\gamma\), \citet{Gaiffas2005DegenerateDesign} obtains the pointwise rate \(n^{-s/(1+2s+\gamma)}\), up to slowly varying factors; the classical positive-density Hölder rate \(n^{-s/(2s+1)}\) from \citet{Stone1982} and the local-polynomial constructions surveyed by \citet{Fan1996} provide the interior reference. In the present boundary regime, the information balance \(n h_n^{2\beta+1}(\delta_n+h_n)^\kappa\asymp1\) recovers the degenerate-design bandwidth when \(\delta_n\) is at the edge scale. Farther into the interior it gives \(h_n\asymp(n\delta_n^\kappa)^{-1/(2\beta+1)}\). The clamp target then multiplies the resulting boundary-regression error by the collapsed mass of order \(\delta_n^{\kappa+1}\), yielding the atom term \(\delta_n^{\kappa+1}h_n^\beta\) alongside the retained-mean \(n^{-1/2}\) term. This mass weighting, absent from ordinary pointwise regression, generates the critical and atom-dominated threshold regimes.  Two antecedents deserve a closer comparison, because they bear directly on the oracle conditions used here and on the sharp-constant question left open below.

The first is adaptation under degenerate design.  \citet{Gaiffas2007Inhomogeneous} studies pointwise estimation from inhomogeneous data when the smoothness is unknown, and obtains the adaptive rate with the usual logarithmic payment: the exponent matches the known-smoothness degenerate-design rate with \(n\) replaced by \(n/\log n\), so adaptation costs a logarithmic factor rather than a change of exponent.  The procedures in the present paper are class-indexed in the oracle sense: the smoothness exponent and radius, the thinning exponent and envelope constants, and the threshold path are fixed before sampling.  Atom weighting changes the adaptation question rather than inheriting it.  The clamp target multiplies the boundary-regression error by the collapsed mass \(\delta_n^{\kappa+1}\), so a logarithmic loss in the regression component is invisible whenever the retained-mean term \(n^{-1/2}\) dominates, and is felt only above the critical threshold scale; moreover the mass itself depends on the thinning exponent, so an adaptive treatment would have to adapt to \(\kappa\) as well as to \(\beta\).  The frontier reported here is the known-class benchmark for that question: it holds with \(\beta\), \(L\), and \(\kappa\) fixed and available to the procedure, which is the reference an adaptive treatment would be measured against.  \Cref{sec:discussion-open-questions} takes up adaptation to unknown \(\beta\), \(L\), and \(\kappa\) as a question for future work.

The second is the modulus of continuity.  The classical account of minimax rates through a modulus is due to \citet{DonohoLiu1991GeometrizingII}, who show that the difficulty of estimating a linear functional over a convex class is governed by the between-class modulus of the functional with respect to a statistical distance, and \citet{CaiLow2004AdaptationTheory} develop the corresponding theory for confidence intervals, where expected length over a subclass is governed by a within-class or between-class modulus rather than by the minimax rate alone.  The honest expected-length lower bound established here is proved by direct two-point and localized-perturbation arguments, and it matches the risk frontier in order.  Relating that frontier to a modulus is a well-posed next step, and the structure of the clamp target says what such a modulus would have to accommodate: the functional is linear in the outcome regression, while its weight \(q_x(\delta)\) is determined by the design, so the appropriate family is a design-indexed one rather than the convex class of a single fixed linear functional.  The sharp-constant question stated in the discussion is exactly this calibration question.

The inference results also build on the theory of honest confidence intervals. \citet{Low1997NonparametricConfidenceIntervals} shows how expected length lower bounds constrain honest adaptation over nonparametric classes, and \citet{ArmstrongKolesar2018OptimalInference,ArmstrongKolesar2020} develop bias-aware interval methods that balance stochastic error and worst-case approximation bias. Those ideas are especially natural here because the clamp target contains a boundary regression term whose bias is controlled by smoothness and whose variance is controlled by the local design mass. The paper establishes matched minimax absolute-risk and uniformly honest expected-length rates for this target, so the estimation and interval scales agree under the same polynomial-thinning conditions.

Finally, the deterministic lower-threshold clamp studied here has a mixed-measure post-intervention law: values already above the threshold remain on their continuous natural scale, while values below the threshold collapse to the boundary point. That perspective explains why the lower-threshold clamp differs from smooth dose-response estimation: the intervention produces an atom whose mass is estimable from treatment frequencies, while its contribution to the mean depends on a boundary response value. The analysis below isolates this atom-regression product and derives the rates for both the Hölder model and the continuity-only model, thereby connecting modified-policy causal estimands to degenerate-design regression and honest nonparametric inference.
